\documentclass[12pt,a4paper]{article}

\usepackage[utf8]{inputenc}
\usepackage[T1]{fontenc}
\usepackage[english]{babel}
\usepackage{geometry}
\usepackage{setspace}
\usepackage{array}
\usepackage{longtable}
\usepackage{booktabs}
\usepackage{ragged2e}
\usepackage{indentfirst}
\usepackage{titlesec}
\usepackage{tocloft}

\geometry{
    left=30mm,
    right=15mm,
    top=20mm,
    bottom=20mm
}

\usepackage[
  unicode,
  pdftitle={Technical task},
  pdfauthor={Artem Davydchuk}
]{hyperref}

\setlength{\parindent}{1.25cm}
\onehalfspacing

\titleformat{\section}
    {\normalfont\bfseries}
    {\thesection.}
    {0.5em}
    {}

\renewcommand{\contentsname}{Contents}

\begin{document}

    \begin{center} \textbf{\Large Technical Specification for the Course Project} \end{center}

    \vspace{0.5cm}

    \tableofcontents

    \newpage

    \section{NAME AND SCOPE OF APPLICATION}

    This technical specification applies to the development of the software application
    \textbf{``System for Dialog Interaction with Large Language Models with Web and Console Interfaces''}.

    The scope of application of the developed software product includes:
    organization of user interaction with a language model,
    management of chat history, launching the application in different operating modes,
    as well as its use for educational, demonstration, and research purposes
    in the fields of information technology, artificial intelligence, and software engineering.

    The software application is intended for local or network use
    on a personal computer and must provide convenient interaction with the system
    both through a browser-based web user interface and through a command-line interface.

    \section{REASONS FOR DEVELOPMENT}

    The basis for development is the assignment for the course project
    within the educational and professional program
    \textbf{``Computer Systems and Networks''}
    in the specialty

    \noindent
    \textbf{123 ``Computer Engineering''}.

    The need for development is due to the growing relevance of systems
    that provide interaction between humans and intelligent software tools,
    in particular with large language models, as well as by the need to create
    a convenient, flexible, and accessible software tool
    that supports multiple modes of operation --- a web interface and a console interface.

    \section{PURPOSE AND OBJECTIVE OF DEVELOPMENT}

    The purpose of this project is to develop a software application for dialog interaction
    with a large language model, which provides:

    \begin{itemize}
        \item operation in two modes: through a browser-based web user interface and through a command-line interface (CLI);
        \item creation of a new chat and selection of an existing chat;
        \item input, transmission, and display of messages within a dialogue;
        \item storage of chat history;
        \item ease of use and clear organization of the interface.
    \end{itemize}

    The objective of the development is to create a universal software tool
    that can be used as a demonstration, educational, or basic applied system
    for working with language models and for further extension of functional capabilities.

    \section{DEVELOPMENT SOURCES}

    The sources for development are:

    \begin{itemize}
        \item scientific and technical literature on software engineering, artificial intelligence, and the construction of dialogue systems;
        \item reference documentation for the Python programming language;
        \item documentation for libraries used to implement web and console interfaces;
        \item Internet publications devoted to the construction of LLM systems, client applications, and user interfaces;
        \item materials on software architecture design and data storage organization.
    \end{itemize}

    \section{TECHNICAL REQUIREMENTS}

    \subsection{Requirements for the Developed Product}

    The developed software application must meet the following requirements:

    \begin{enumerate}
        \item The program must be launched from the main file \texttt{app.py}.
        \item Two modes of operation must be implemented:
        \begin{itemize}
            \item standard launch --- launching the web user interface accessible through a browser;
            \item launch with the \texttt{-t} parameter --- switching to command-line mode.
        \end{itemize}
        \item In web interface mode, the following must be provided:
        \begin{itemize}
            \item display of the chat list;
            \item creation of a new chat;
            \item opening an existing chat;
            \item display of message history;
            \item input of a new message and receipt of a response.
        \end{itemize}
        \item In console mode, the following must be provided:
        \begin{itemize}
            \item creation of a new chat;
            \item selection of an existing chat;
            \item exchange of messages in text mode;
            \item correct termination of the session.
        \end{itemize}
        \item Chat history storage in files or local data storage must be implemented.
        \item The program interface must be logically structured, understandable, and convenient for the user.
        \item The program must correctly handle input errors, startup failures, and invalid command-line parameters.
        \item The software architecture must allow for further extension of functionality.
    \end{enumerate}

    \subsection{Software Requirements}

    The following are required for the functioning of the software application:

    \begin{itemize}
        \item operating system: Windows 10/11, Linux, or another modern operating system with Python support;
        \item Python interpreter version 3.10 or higher;
        \item a modern web browser with support for HTML5, CSS, and JavaScript;
        \item libraries and modules required to implement the web user interface, program logic, and chat history storage;
        \item software development and debugging tools (if necessary).
    \end{itemize}

    \subsection{Hardware Requirements}

    For operation of the software application, a personal computer is required
    with the following minimum characteristics:

    \begin{itemize}
        \item x86\_64 architecture processor;
        \item RAM --- at least 8 GB;
        \item free disk space --- at least 1 GB;
        \item monitor, keyboard, and input device;
        \item if necessary --- Internet connection for working with external services or models.
    \end{itemize}

    \section{DEVELOPMENT STAGES}

    \begin{longtable}{|>{\centering\arraybackslash}m{1.5cm}|>{\RaggedRight\arraybackslash}m{10.5cm}|>{\centering\arraybackslash}m{3cm}|}
        \hline
        \textbf{No.} & \textbf{Development Stage} & \textbf{Date} \\
        \hline
        6.1 & Receiving the topic and assignment & 15.02.2026\\
        \hline
        6.2 & Selection and study of literature on dialogue systems, LLMs, and web user interfaces & 15.03.2026\\
        \hline
        6.3 & Preparation of the technical specification & 15.03.2026\\
        \hline
        6.4 & Analysis of the subject area and review of existing approaches to building dialogue applications & \\
        \hline
        6.5 & Designing the structure of the software application, web interface and CLI modes, as well as the chat storage mechanism & \\
        \hline
        6.6 & Development of the software application & \\
        \hline
        6.7 & Testing of the software application & \\
        \hline
        6.8 & Preparation of the explanatory note & \\
        \hline
        6.9 & Submission of the course project (paper) for review & \\
        \hline
        6.10 & Defense of the course project (paper) & \\
        \hline
    \end{longtable}

\end{document}